\documentclass[a4paper,11pt]{article}
\usepackage{exam-style-842-unit}
\examinehead{842 信息化概论 · 第 2 单元练习册}{练习二 参考答案与解析}
\begin{document}

\unitanswerscover{842 信息化概论 · 单元练习册}{第 2 单元参考答案}{练习二（信息系统与数据资源管理）的答案与解析；每道判断题均给出改正后的正确表述。}

\answerpart{一、单项选择题（10 题，每题 3 分，共 30 分）}

\q D。\begin{analysis}信息系统的五个基本功能（基本部分）是信息的输入、处理、输出、存储和控制（大纲「一、信息技术」第 3 条：信息系统组成）。资金的筹措属于财务活动，不是信息系统的功能。A、B、C 分别对应输入、处理与存储功能。\end{analysis}

\q B。\begin{analysis}信息系统是由人、硬件、软件、数据和处理规程（流程）五个要素组成的人机系统，其中处理规程规定信息处理与系统运行的规则和步骤（大纲「一、信息技术」第 3 条；通行信息系统理论）。A 网络是基础设施的一部分，C 机房是运行环境，D 报表是输出形式，均不是与前三者并列的组成要素。\end{analysis}

\q B。\begin{analysis}信息资源管理的核心是把信息看作与物质、能源同等重要的战略资源，对其进行合理开发、配置和有效利用，以提高组织效益和决策水平（大纲「一、信息技术」第 3 条；通行信息资源管理理论）。A 数字化只是手段；C 只强调设备投入、D 只强调建库，都属于以技术代替管理的片面理解。\end{analysis}

\q A。\begin{analysis}信息采集是信息资源管理的第一个环节，其基本要求是全面、准确、及时：全面指采集范围覆盖管理需要，准确指数据真实可靠，及时指信息在有效期内进入系统（大纲「信息资源获取、交互与资源共享」；通行信息资源管理理论）。B 只讲数量和成本，C 把采集等同于保密，D 与采集要求相反，均不全面。\end{analysis}

\q A。\begin{analysis}事务处理系统（TPS）面向组织日常业务，处理订单、账务、库存等大量结构化事务数据，是信息系统的基础层次。B 管理信息系统在 TPS 基础上加工管理信息；C 决策支持系统面向半结构化决策；D 专家系统以知识库和推理机模拟专家解决问题，都不是处理日常事务数据的系统。\end{analysis}

\q B。\begin{analysis}管理信息系统（MIS）在事务处理系统的基础上，对管理信息进行加工处理，为各级管理者提供报表、查询和管理控制支持。A 事务处理系统只处理日常事务；C 决策支持系统强调模型与半结构化决策；D 办公自动化系统主要处理办公事务与文档流转，均与题干描述不符。\end{analysis}

\q C。\begin{analysis}决策支持系统（DSS）以数据、模型和方法为手段，通过人机交互辅助管理者解决半结构化、非结构化决策问题。A 事务处理系统面向结构化日常事务；B 管理信息系统面向例行管理信息与报表；D 数据库系统只是数据管理的支撑平台，不直接提供决策模型。\end{analysis}

\q A。\begin{analysis}信息安全的基本属性是机密性、完整性、可用性（CIA）：机密性保证信息不被非授权者获取，完整性保证信息不被篡改或破坏，可用性保证授权者能正常使用信息与系统（教材 6.3.3，第 110—112 页；大纲「信息的安全与保护」）。B 是信息共享特性的延伸表述；C 是信息采集的要求；D 是系统选型的考虑因素，均不是信息安全的基本属性。\end{analysis}

\q B。\begin{analysis}数据备份是为防止数据丢失或损坏，把数据复制到其他存储介质上保存的保护措施，是数据安全与系统恢复的基础工作。A 加密防止数据被窃取后泄露内容；C 压缩用于节省存储与传输空间；D 索引用于加快检索，均不能替代备份的防丢失作用。\end{analysis}

\q B。\begin{analysis}信息具有可共享性：同一信息可以被多个用户同时使用而不会损耗，共享范围越大信息价值越高（教材 1.1.2，第 3 页）。A 可存储性指信息可记录保存；C 可压缩性指可去除冗余；D 可度量性指信息量可用熵等指标度量，均与“多人同时使用而不损耗”不符。\end{analysis}

\answerpart{二、多项选择题（3 题，每题 4 分，共 12 分）}

\q ABC。\begin{analysis}信息系统由人、硬件、软件、数据和处理规程五个要素组成。A 人、B 硬件与软件、C 数据都是组成要素；D 错误：资金是信息化建设的投入条件，但不是信息系统的组成要素，把资金列为组成要素会混淆“建设资源”与“系统构成”。（大纲「一、信息技术」第 3 条；通行信息系统理论）\end{analysis}

\q ABC。\begin{analysis}信息资源管理是对信息、信息人员、信息技术与设施等要素进行计划、组织、协调与控制，内容包括信息资源的开发与利用、信息基础设施与信息系统建设、信息安全与保密管理等。D 错误：信息资源管理既包括电子资源也包括纸质资源，但绝不“只管理纸质档案”，恰恰相反，其重点是用信息技术管理各类信息资源。\end{analysis}

\q ABC。\begin{analysis}信息安全的基本属性是机密性、完整性和可用性（教材 6.3.3）；广义上还包括可控性。D 错误：“随意性”与信息安全的要求相反，信息访问必须经过授权和审计，不能随意。\end{analysis}

\answerpart{三、判断题（6 题，每题 2 分，共 12 分）}

\q ×。改正：信息系统是由人、硬件、软件、数据和处理规程组成的人机系统，计算机硬件只是其中的组成部分之一。\begin{analysis}信息系统不仅有硬件，还包括软件、数据、人员和规程；只强调硬件会忽略人在系统中的主体作用和数据的核心地位。（大纲「一、信息技术」第 3 条）\end{analysis}

\q ×。改正：信息具有可共享性，同一信息可被多个用户同时使用而不损耗，共享范围越大信息价值越高。\begin{analysis}与物质、能源不同，信息不因使用而消耗；信息共享是信息资源开发利用的重要目标，也是信息区别于物质资源的基本特性。（教材 1.1.2，第 3 页）\end{analysis}

\q √。\begin{analysis}数据备份把数据复制到其他存储介质上保存，当原数据因故障、误操作或灾害丢失时可以由备份恢复，是防止数据丢失和损坏的有效措施，也是数据安全与灾难恢复的基础。\end{analysis}

\q ×。改正：信息系统的输出形式多种多样，可以是屏幕显示、打印报表、数据文件、控制信号等，并不限于纸质报表。\begin{analysis}输出是信息系统把处理结果显示给用户或传递给其他系统的环节，输出形式和载体随应用需要而定；把输出等同于纸质报表会窄化信息系统的功能。\end{analysis}

\q √。\begin{analysis}信息采集要求全面、准确、及时：全面保证管理所需信息不遗漏，准确保证数据真实可靠，及时保证信息在有效期内进入系统，三者共同决定后续加工与决策的质量。\end{analysis}

\q √。\begin{analysis}决策支持系统以数据、模型和方法为手段，通过人机交互辅助管理者进行半结构化、非结构化决策，强调支持而不是代替决策者判断。\end{analysis}

\answerpart{四、填空题（6 空，每空 1 分，共 6 分）}

\q （1）控制。\begin{analysis}信息系统的五个基本功能是输入、处理、输出、存储和控制，控制功能对系统运行进行监测和调节，保证系统按预定目标工作。\end{analysis}

\q （2）处理规程。\begin{analysis}信息系统的五个组成要素是人、硬件、软件、数据和处理规程（流程）。处理规程指信息处理与系统运行所遵循的规则、步骤和制度。\end{analysis}

\q （3）共享。\begin{analysis}信息资源的开发利用包括信息的采集、加工、存储、检索、传输和共享利用等环节，共享是信息资源价值实现的关键环节。（大纲「信息资源获取、交互与资源共享」）\end{analysis}

\q （4）可用性。\begin{analysis}信息安全的基本属性是机密性、完整性和可用性。可用性指保证信息与信息系统能被授权者正常使用。\end{analysis}

\q （5）事务处理系统。\begin{analysis}事务处理系统（TPS）面向组织日常业务，处理大量结构化事务数据，是管理信息系统和决策支持系统的数据基础。\end{analysis}

\q （6）数据备份。\begin{analysis}数据备份是把数据复制到其他存储介质上保存，以便在原数据丢失或损坏时恢复，是数据安全保护的基础措施。\end{analysis}

\answerpart{五、简答题（4 题，每题 6 分，共 24 分）}

\q 参考答案要点：
（1）信息系统是由人、硬件、软件、数据和处理规程组成，对信息进行采集、加工、存储、传输和利用的人机系统，其目的是支持组织的管理、决策和业务活动。
（2）组成要素：人（系统建设者与使用者）、硬件（计算机与网络等设备）、软件（系统软件与应用软件）、数据（系统的处理对象与核心资源）、处理规程（信息处理与系统运行的规则、步骤和制度）。
（3）五个基本功能：输入（采集并录入信息）、处理（对信息进行分类、计算、分析等加工）、输出（以报表、屏幕、文件、控制信号等形式提供结果）、存储（长期保存数据与程序）、控制（监测和调节系统运行，保证系统正确高效工作）。
\begin{analysis}得分点：信息系统定义 2 分，五个组成要素 2 分，五个基本功能 2 分。常见失分：把资金、网络单列为组成要素；只答“输入—处理—输出”而漏掉存储与控制。大纲「一、信息技术」第 3 条要求掌握信息系统组成；教材无信息系统专节，本答案按通行信息系统理论作答（见本单元 README）。\end{analysis}

\q 参考答案要点：
（1）信息资源管理的含义：把信息视为与物质、能源同等重要的战略资源，对信息活动中的信息、人员、技术、设备和资金等要素进行计划、组织、协调与控制，实现信息资源的合理开发与有效利用。
（2）主要目标：以尽可能低的成本，及时、准确、完整地向用户提供所需信息，提高组织的管理水平和决策能力；实现信息资源共享，避免重复采集和重复建设；保障信息安全。
（3）主要内容：信息资源的开发与利用（采集、加工、存储、检索、传输、共享）；信息基础设施与信息系统的建设与运行管理；信息资源的标准化与规范化；信息安全与保密管理；信息人员与信息机构的组织管理；信息政策、法规与制度建设。
\begin{analysis}得分点：含义 2 分，目标 2 分，内容 2 分。常见失分：把信息资源管理等同于数据库管理（对象过窄）；只谈技术手段而漏掉人员、制度与政策。依大纲「信息资源管理技术（信息资源获取、交互与资源共享等），信息的安全与保护」及通行信息资源管理理论作答。\end{analysis}

\q 参考答案要点：
（1）信息采集：通过传感、调查、统计、录入等手段获取原始信息，要求全面、准确、及时。
（2）信息加工：对原始信息进行分类、排序、计算、分析、综合等处理，使之成为对使用者有用的信息，要求真实、适用、及时。
（3）信息存储：把信息以一定格式记录在相应载体上长期保存，以便反复使用，要求安全、可靠、便于检索。
（4）信息检索：从信息集合中查找所需信息，有目录检索、关键词检索、全文检索等方式，要求快速、准确、查全率高。
（5）信息传输：信息从信源经信道传递到信宿，要求快速、可靠、安全。
（6）信息共享：信息可被多个用户同时使用而不损耗，要求打破部门壁垒、统一标准、明确权限，在共享的同时保护信息安全。
\begin{analysis}得分点：六个环节各 1 分。常见失分：漏掉“检索”或“传输”环节；只写环节名称而不写基本要求；把“共享”与“安全”对立起来。依大纲「信息资源获取、交互与资源共享」作答。\end{analysis}

\q 参考答案要点：
（1）信息安全的基本属性：机密性（保证信息不被非授权者获取，主要靠密码技术）、完整性（保证信息不被篡改或破坏）、可用性（保证授权者能正常使用信息与系统）；广义上还包括可控性。
（2）常用数据保护措施：
（1）技术措施：加密、身份认证与授权、访问控制、数据备份与恢复、审计与日志、防病毒与防火墙、入侵检测等；
（2）管理措施：制定数据安全管理制度和操作规程，明确数据分级分类与使用权限，落实安全责任制，定期开展安全检查和应急演练；
（3）法律与制度措施：遵守网络安全与数据保护相关法律法规，保护用户隐私和个人信息；
（4）物理与环境措施：机房环境安全、电磁防护、介质管理等。
（3）三属性之间存在制约关系：过强的机密性要求可能影响可用性，必须在安全与效率之间取得平衡，按数据重要程度分级保护。
\begin{analysis}得分点：三个基本属性各 1 分共 3 分，保护措施答出技术与管理两方面共 2 分，指出安全与效率的平衡 1 分。常见失分：把“信息安全”只理解为防病毒；只答技术措施而忽略管理制度。依教材 6.3.3（第 110—112 页）与大纲「信息的安全与保护」作答。\end{analysis}

\answerpart{六、综合题（论述，1 题，共 16 分）}

\q 参考答案要点：
（1）组成要素与基本功能（4 分）：该平台作为信息系统，由人（农业农村局管理人员、农技人员、运维人员）、硬件（服务器、网络、终端、传感器）、软件（数据库管理系统、平台应用软件、分析工具）、数据（土壤、气象、市场行情、种植户等数据）和处理规程（数据采集规范、共享审批流程、安全管理制度）组成。平台的功能包括：输入——采集土壤、气象、市场等数据并录入；处理——清洗、汇总、分析数据，生成农情与市场分析；输出——向乡镇、合作社、企业提供查询、报表、预警信息；存储——集中存储各类数据；控制——通过权限管理、日志审计和运行监控保证平台安全稳定运行。
（2）目标与主要环节（4 分）：信息资源管理的目标是把信息作为战略资源合理开发、有效利用，以较低成本及时、准确、完整地提供信息，支持农业生产与政府决策。该平台的主要环节包括：信息采集（布设监测点、调查统计、系统对接）、信息加工（数据清洗、统计分析、模型计算）、信息存储（集中数据库与备份）、信息检索（按作物、区域、时间查询）、信息传输（通过政务网、互联网向用户推送）、信息共享（向乡镇、合作社、企业分级开放），并全程贯穿安全保护。
（3）共享与安全（4 分）：开放共享必须以安全为前提：明确数据分类分级，对涉及农户隐私和商业秘密的数据进行脱敏或限制开放；实行身份认证与基于角色的访问控制，不同用户授予不同权限；数据传输和存储采用加密技术，保证机密性和完整性；建立数据备份与灾难恢复机制，保证可用性；通过日志审计和安全监测发现异常访问；制定数据共享与安全管理制度，明确责任主体和使用规范，做到“共享有度、使用留痕”。
（4）信息孤岛问题与对策（4 分）：建设中可能出现各部门、各乡镇系统独立建设、数据标准不一、接口不通，形成信息孤岛，导致数据重复采集、口径矛盾、共享困难。对策包括：统一数据标准与编码规范；建设统一的数据库与共享交换平台，提供标准接口；建立跨部门数据共享协调机制和权责清单；开展数据治理，做好数据质量核查与更新维护；在制度上明确数据“一数一源”和共享责任，避免重复建设。
\begin{analysis}得分点：（1）五要素与五功能 4 分；（2）信息资源管理目标与采集、加工、存储、检索、传输、共享六环节 4 分；（3）安全属性与分级分类、权限控制、加密备份、审计制度等 4 分；（4）信息孤岛成因与标准、平台、机制、治理等对策 4 分。常见失分：把平台功能只写成“存数据、查数据”；安全部分只写“加密”而漏掉管理与制度；信息孤岛对策只喊口号而无具体措施。大纲第一板块第 3 条与教材 6.3.3 为本题主要依据，信息系统组成按通行理论补充（见本单元 README）。\end{analysis}

\end{document}
